\section{Database Design}

\subsection{Business Requirement Analysis}

\subsubsection{Business Purpose}

The system manages the booking and usage of shared campus spaces for the School of Computer Science, replacing a manual spreadsheet/email process with a structured database that handles booking requests, approvals, check-in/check-out, maintenance tracking, and historical reporting.

\subsubsection{Actors}

\begin{description}
	\item[Student] Submits booking requests for student activities; views own bookings.
	\item[Lecturer] Submits booking requests for lectures and examinations.
	\item[Teaching Assistant] Submits booking requests for workshops and tutorials.
	\item[Facility Staff] Approves/rejects bookings, performs check-in/check-out, manages maintenance.
	\item[Department Administrator] Views departmental booking history.
	\item[Facility Manager] Oversees all spaces, bookings, maintenance, and facility statuses.
\end{description}

\subsubsection{Core Entities}

% Helper: monospace name with a breakpoint allowed after it
\newcommand{\col}[1]{\texttt{#1}\allowbreak}

\begin{description}[style=nextline]
	\begingroup\raggedright
	\item[\texttt{USER}]
	      \col{user\_id} (PK), \col{full\_name}, \col{email}, \col{phone}, \col{role}, \col{department}, \col{account\_status}.

	\item[\texttt{SPACE}]
	      \col{space\_code} (PK), \col{space\_name}, \col{space\_type}, \col{building}, \col{floor}, \col{room\_number}, \col{capacity}, \col{current\_status}, \col{usage\_policy}.

	\item[\texttt{FACILITY}]
	      \col{facility\_id} (PK), \col{space\_code} (FK), \col{facility\_name}, \col{description}.

	\item[\texttt{BOOKING\_REQUEST}]
	      \col{booking\_id} (PK), \col{user\_id} (FK), \col{space\_code} (FK), \col{requested\_start\_time}, \col{requested\_end\_time}, \col{purpose}, \col{expected\_participants}, \col{booking\_type}, \col{status}.

	\item[\texttt{BOOKING\_APPROVAL}]
	      \col{approval\_id} (PK), \col{booking\_id} (FK, UNIQUE), \col{decided\_by\_user\_id} (FK), \col{decision\_time}, \col{decision\_note}, \col{rejection\_reason}.

	\item[\texttt{USAGE\_SESSION}]
	      \col{session\_id} (PK), \col{booking\_id} (FK, UNIQUE), \col{actual\_start\_time}, \col{actual\_end\_time}, \col{checked\_in\_by\_user\_id} (FK), \col{completed\_by\_user\_id} (FK), \col{initial\_condition}, \\ \col{final\_condition}, \col{usage\_notes}.

	\item[\texttt{MAINTENANCE\_RECORD}]
	      \col{maintenance\_id} (PK), \col{space\_code} (FK), \col{reporter\_user\_id} (FK), \\ \col{assigned\_staff\_user\_id} (FK), \col{problem\_description}, \col{start\_time}, \\ \col{completion\_time}, \col{status}, \col{result\_note}.
	      \endgroup
\end{description}

\subsubsection{Business Rules}

\begin{enumerate}
	\item No overlapping approved bookings for the same space.
	\item Spaces with status \texttt{under\_maintenance}, \texttt{temporarily\_closed}, or \texttt{retired} cannot be booked.
	\item Spaces with active maintenance records cannot be booked.
	\item \texttt{requested\_end\_time} must be greater than \texttt{requested\_start\_time}.
	\item \texttt{BOOKING\_APPROVAL} is optional, at most one per \texttt{BOOKING\_REQUEST} (1:0..1).
	\item \texttt{USAGE\_SESSION} is optional, at most one per \texttt{BOOKING\_REQUEST} (1:0..1).
	\item Rejected bookings must preserve \texttt{rejection\_reason}.
	\item Check-in must record \texttt{actual\_start\_time}, \texttt{checked\_in\_by\_user\_id}, \\ \texttt{initial\_condition}.
	\item Completion must record \texttt{actual\_end\_time}, \texttt{completed\_by\_user\_id}, \texttt{final\_condition}, \texttt{usage\_notes}.
	\item Historical data must be preserved (no hard deletes).
\end{enumerate}

\subsubsection{Requirement Summary}

\begin{description}
	\item[User management] \texttt{USER} with role and account status.
	\item[Space inventory] \texttt{SPACE} with type, capacity, and status.
	\item[Facility tracking] \texttt{FACILITY} linked to \texttt{SPACE} via FK.
	\item[Booking workflow] \texttt{BOOKING\_REQUEST} with status lifecycle.
	\item[Approval management] \texttt{BOOKING\_APPROVAL} (1:0..1 to request).
	\item[Usage tracking] \texttt{USAGE\_SESSION} for check-in/check-out.
	\item[Maintenance management] \texttt{MAINTENANCE\_RECORD} from report to resolution.
	\item[Conflict prevention] Triggers/application logic on approved bookings.
	\item[History preservation] Soft-delete policy across all relations.
\end{description}

\subsection{Conceptual Database Design}

\subsubsection{Entity Relationship Diagram}

The Entity Relationship Diagram (ERD) models the conceptual structure of the Campus Space Management System. It consists of seven entities, together with their attributes, relationships, cardinalities, and participation constraints. The complete Mermaid source is provided in the project deliverables (\texttt{output/02-erd-design-G08.md}), while Figure~\ref{fig:erd} presents the rendered ERD.

\begin{figure}[H]
    \centering
    \includegraphics[width=\textwidth]{src/diagram.png}
    \caption{Entity Relationship Diagram of the Campus Space Management System}
    \label{fig:erd}
\end{figure}

\subsubsection{Design Decisions}

\begin{itemize}
	\item \textbf{Entity identification:} Seven entities map one-to-one to distinct real-world concepts, no merging or over-splitting.
	\item \textbf{Relationship design:} \texttt{USER} participates in multiple relationships with distinct semantic roles (requester, approver, check-in staff, reporter, assigned staff), modeled via separate foreign keys.
	\item \textbf{Cardinalities:} \texttt{BOOKING\_APPROVAL} and \texttt{USAGE\_SESSION} use 1:0..1 (not 1:N), reflecting at-most-one approval/session per booking.
	\item \textbf{Participation:} Total participation for dependent entities (e.g., every \texttt{FACILITY} must belong to a \texttt{SPACE}); partial participation where entities may exist independently.
\end{itemize}

\subsubsection{Relationship Summary}

\begin{itemize}
	\item Space contains Facility \dotfill 1 : N
	\item User submits Booking\_Request \dotfill 1 : N
	\item Space receives Booking\_Request \dotfill 1 : N
	\item Booking\_Request has Booking\_Approval \dotfill 1 : 0..1
	\item User performs Booking\_Approval \dotfill 1 : N
	\item Booking\_Request creates Usage\_Session \dotfill 1 : 0..1
	\item User checks in Usage\_Session \dotfill 1 : N
	\item User completes Usage\_Session \dotfill 1 : N
	\item Space has Maintenance\_Record \dotfill 1 : N
	\item User reports Maintenance\_Record \dotfill 1 : N
	\item User assigned to Maintenance\_Record \dotfill 1 : N
\end{itemize}


\subsection{Logical Database Design}

\subsubsection{Relational Schema}

The ERD was converted into seven relations. The full schema definitions with data types are provided in the project deliverables (\texttt{output/03-logical-design-G08.md}).

\subsubsection{Relation Summary}

\begin{description}[style=nextline]
	\begingroup\raggedright
	\item[\texttt{USERS}]
	      PK: \col{user\_id}.
	      Candidate key: \col{email}.

	\item[\texttt{SPACES}]
	      PK: \col{space\_code}.
	      Candidate key: (\col{building}, \col{room\_number}).

	\item[\texttt{FACILITY}]
	      PK: \col{facility\_id}.
	      FK: \col{space\_code}.

	\item[\texttt{BOOKING\_REQUEST}]
	      PK: \col{booking\_id}.
	      FKs: \col{user\_id}, \col{space\_code}.

	\item[\texttt{BOOKING\_APPROVAL}]
	      PK: \col{approval\_id}.
	      FKs: \col{booking\_id} (UNIQUE), \col{decided\_by\_user\_id}.
	      Candidate key: \col{booking\_id}.

	\item[\texttt{USAGE\_SESSION}]
	      PK: \col{session\_id}.
	      FKs: \col{booking\_id} (UNIQUE), \col{checked\_in\_by\_user\_id}, \col{completed\_by\_user\_id}.
	      Candidate key: \col{booking\_id}.

	\item[\texttt{MAINTENANCE\_RECORD}]
	      PK: \col{maintenance\_id}.
	      FKs: \col{space\_code}, \col{reporter\_user\_id}, \\ \col{assigned\_staff\_user\_id}.
	      \endgroup
\end{description}

All eleven foreign keys implement the ERD relationships. The UNIQUE constraint on \texttt{booking\_id} in \texttt{BOOKING\_APPROVAL} and \texttt{USAGE\_SESSION} enforces the 1:0..1 cardinality.

\subsubsection{Normalization}

All relations satisfy Third Normal Form (3NF):

\begin{itemize}
	\item \textbf{1NF:} All attributes are atomic.
	\item \textbf{2NF:} No partial dependencies (all PKs are single-attribute).
	\item \textbf{3NF:} No transitive dependencies among non-key attributes.
\end{itemize}

\subsection{Database Validation}

\subsubsection{Requirement Traceability}

Every business requirement maps to a database component:
user management to \texttt{USERS},
space inventory to \texttt{SPACES},
facility tracking to \texttt{FACILITY},
booking workflow to \texttt{BOOKING\-\_REQUEST},
approval decisions to \texttt{BOOKING\-\_APPROVAL},
usage tracking to \texttt{USAGE\-\_SESSION},
and maintenance management to \texttt{MAINTENANCE\-\_RECORD}.

\subsubsection{Constraint Analysis}

\begin{table}[H]
	\centering
	\small
	\begin{tabularx}{\textwidth}{|>{\raggedright\arraybackslash}X|>{\raggedright\arraybackslash}X|>{\raggedright\arraybackslash}X|}
		\hline
		\textbf{Business Rule}            & \textbf{Schema Constraint}        & \textbf{Additional Enforcement} \\
		\hline
		No overlapping bookings           & ---                               & Trigger / application logic     \\
		\hline
		Unavailable spaces not bookable   & CHECK on \texttt{current\_status} & Application status check        \\
		\hline
		Active maintenance blocks booking & ---                               & Application logic               \\
		\hline
		End time > start time             & CHECK constraint                  & ---                             \\
		\hline
		At most one approval per booking  & UNIQUE on \texttt{booking\_id}    & ---                             \\
		\hline
		At most one session per booking   & UNIQUE on \texttt{booking\_id}    & ---                             \\
		\hline
		Rejection reason preserved        & ---                               & Application logic               \\
		\hline
		Valid user roles                  & CHECK on \texttt{role}            & ---                             \\
		\hline
		No hard deletes                   & ---                               & Application policy              \\
		\hline
	\end{tabularx}
\end{table}

\subsubsection{Design Limitations}

The following constraints cannot be expressed as single-table CHECK constraints and require triggers or application logic:

\begin{itemize}
	\item \textbf{Overlapping booking prevention:} Requires querying existing approved bookings for time-range overlap detection.
	\item \textbf{Maintenance-based booking blocks:} Requires cross-table join between \texttt{BOOKING\_REQUEST} and \texttt{MAINTENANCE\_RECORD}.
	\item \textbf{Conditional NOT NULL:} Fields like \texttt{rejection\_reason} and completion fields are conditionally required based on booking status.
\end{itemize}
